\section{Apparatus and Study Setting}
\label{sec:study_design_apparatus}

The apparatus consists of the physical rowing ergometer, the VR headset, the tracking and calibration setup, and the Unity-based rowing prototype described in Chapter~\ref{chap:system}. The same physical rowing task and the same basic prototype environment will be used in both conditions. This reduces the risk that differences in ratings are caused by changes in exercise task, exertion context, or study setting rather than by the representation mode.

\drafttodo{TODO:CONFIRM exact ergometer model, headset, tracking devices, computer or standalone setup, software version, and room setup.}

The participant interacts with the real ergometer while wearing the headset. The handle relation is tracked through the participant's hands, while the remaining rowing-body motion is inferred from the aligned rowing machine, calibrated setup geometry, and participant body measurements. This makes calibration part of the experimental apparatus rather than a neutral preparation step: if the machine, body scale, or virtual rowing frame are poorly aligned, the resulting ratings may reflect setup error rather than the intended contrast between ergometer-congruent and boat-realistic representation.

\drafttodo{TODO:CONFIRM exact hand tracking method, body-measurement inputs, calibration steps, approximate calibration duration, and any measured or estimated end-to-end latency.}

Latency and responsiveness will be documented where possible because delayed closed-loop avatar feedback can influence agency, ownership, simultaneity judgments, and motor performance \parencite{waltemateEtAl2016latency}.
% SOURCE-CHECK: MD Papers/waltemate_2016_latency_perceptual_judgments_motor_performance_vr_markdown_bundle/waltemate_2016_latency_perceptual_judgments_motor_performance_vr_thesis_notes.md
